\documentclass[11pt,a4paper]{article}
\usepackage[utf8]{inputenc}
\usepackage[T1]{fontenc}
\usepackage[french]{babel}
\usepackage{amsmath,amssymb,amsthm,mathtools}
\usepackage{booktabs,array,enumitem}
\usepackage[a4paper,margin=2.5cm]{geometry}
\usepackage[hidelinks]{hyperref}
\usepackage[table]{xcolor}
\definecolor{plotkin}{RGB}{255,200,120}   % borne atteinte, cas non trivial
\definecolor{trivial}{RGB}{225,225,225}   % borne atteinte, cas trivial r=m-1


% ---------- macros ----------
\newcommand{\F}{\mathbb{F}_2}
\newcommand{\wt}{\mathrm{wt}}
\newcommand{\Nl}{\mathrm{NL}}
\newcommand{\NQ}{\mathrm{NQ}}
\newcommand{\NC}{\mathrm{NC}}
\newcommand{\dist}{d}
\newcommand{\Bc}{\mathcal{B}}
\def\bst(#1,#2,#3){\mathcal{B}_{#1}^{#2}(#3)}
\def\clst(#1,#2,#3){\tilde{\mathcal{B}}_{#1}^{#2}(#3)}
\def\NL(#1,#2){ {\rm NL}_{#1}(#2) }
\def\AGL(#1,#2){ {\rm AGL}(#1,#2) }
\def\agl#1{ {\mathfrak #1 }}
\def\rmc(#1,#2){ \textsc {rm}(#1,#2) }
\def\boole(#1){\mathcal{B}(#1)}
\def\fd(#1){\mathcal{F}_2^{#1}}

% ---------- theorem environments ----------
\theoremstyle{plain}
\newtheorem{theorem}{Théorème}[section]
\newtheorem{proposition}[theorem]{Proposition}
\newtheorem{result}[theorem]{Numerical fact}
\newtheorem{lemma}[theorem]{Lemme}
\newtheorem{coro}[theorem]{Corollaire}
\theoremstyle{definition}
\newtheorem{definition}[theorem]{Définition}
\newtheorem{hypothesis}[theorem]{Hypothèse}
\theoremstyle{remark}
\newtheorem{remark}[theorem]{Remarque}

\title{majoration de  $\rmc(2,8)$ et $\rmc(3,8)$}
\author{Valérie Gillot, Philippe Langevin, Jean-Baptiste Nguyen\\
{\small IMATH, Université de Toulon}\\[4pt]
{\small Rapport de travail, rédigé avec l'assistance de Claude (Anthropic)}}
\date{Automne 2026}

\begin{document}
\maketitle

\begin{abstract}
Nous montrons comment obtenir et utiliser la classification de 
$\bst(3,5,7)$ pour obtenir deux nouvelles majorations des codes
de Reed-Muller d'ordre 2 et 3 en 8 variables.
\end{abstract}

\section{Introduction}


Le rayon de recouvrement $\rho(r, m)$ de $\rmc(r,m)$ est le maximum, sur les fonctions
$f\in\boole(m)$, de la distance $\dist(f,\rmc(r,m))$, appelée non-linéarité d'ordre $r$ de $f$,
nous la noton $\NL(r,f)$. Tous les rayons de recouvrements sont connus pour $m\leq 7$,
signalons que  la détermination de $\rho(1,9)$ vient d'être faite après 40 ans d'efforts,
alors les deux premiers cas ouverts concernent les codes de Reed-Muller d'ordre 2 et
et 3 en 8 variables :
$$
	88 \leq \rho(2,8) \leq 96, \qquad  56\leq \rho(3,8) \leq 60.
$$


\begin{table}
        \center
        \caption{$\rho(r,m)$, inclut 10 cas non trivaux de Plotkin-extrèmale.}
        \label{tab:tab_rho}
        \begin{tabular}{c|ccccccccc}
                $r\backslash m$ & 3 & 4 & 5 & 6 & 7 & 8 & 9 & 10 & 11 \\
                \hline
                1 & 2 & \cellcolor{plotkin}6 & 12 & \cellcolor{plotkin}28 & 56 & \cellcolor{plotkin}120 & 243 & 496 & 996 \\
                2 & \cellcolor{trivial}1 & 2 & 6 & \cellcolor{plotkin}18 & 40 & 88-96 & 196-216 & 400-460 & 856-956 \\
                3 & 0 & \cellcolor{trivial}1 & 2 & \cellcolor{plotkin}8 & 20 & 56-60 & 111-156 & 194-372 & 454-832 \\
                4 &  & 0 & \cellcolor{trivial}1 & 2 & 8 & 26 & 58-86 & $\leq$ 242 & $\leq$ 614 \\
                5 &  &  & 0 & \cellcolor{trivial}1 & 2 & \cellcolor{plotkin}10 & 28-36 & $\leq$ 122 & $\leq$ 364 \\
                6 &  &  &  & 0 & \cellcolor{trivial}1 & 2 & 10 & $\leq$ 46 & $\leq$ 168 \\
                7 &  &  &  &  & 0 & \cellcolor{trivial}1 & 2 & \cellcolor{plotkin}12 & $\leq$ 58 \\
                8 &  &  &  &  &  & 0 & \cellcolor{trivial}1 & 2 & 12 \\
                \hline
        \end{tabular}
\end{table}

Les majorations sont des conséquences d'inégalité de Plotkin
$$\rho(r,m) \leq \rho(r,m-1) + \rho(r-1,m-1).$$


Dans cette note, nous nous interessons à l'existence des
fonctions Plotkin-extrèmale dont la  non linéarité d'ordre $r$ 
atteignent la majoration de Plotkin i.e. $\NL(r,m) = \rho(r,m-1) + \rho(r-1,m-1)$.
Une fonctionn de $\boole(m)$  s'écrit $f=g+x_m h$ où les
fonction $g$ et $h$ sont des fonctions de $m$ variables, et
la table de vérité de $(f)$ est la concaténation $( f_0 || f_1 )$
des restrictions  $f_0:= h$, $f_1:=g+h$.

Avec ces notations, pour une meilleure approximation de $f_0$
$$
	\wt( f + q  ) = \wt( f_0 +q_0 ) + \wt( f_1 +q_1) \leq \NL(r,m-1) + \wt( f_1 +q_1) 
	\leq \NL(r,m-1) + \NL(r-1, f_1 + q_1 ) 
$$

Soit $g\in\mathrm{rm}(r-1,m-1)$ telle que $d(f_1+q_0,g)=NL_{r-1}(f_1+q_0)$.
Par la construction $(u\,\|\,u+v)$, la fonction $q=(q_0\,\|\,q_0+g)$ appartient à
$\mathrm{rm}(r,m)$, et
\[
d(f,q)=d(f_0,q_0)+d(f_1,q_0+g)=NL_r(f_0)+NL_{r-1}(f_1+q_0).
\]
D'où la première inégalité, la seconde résulte de la définition de $\rho(r-1,m-1)$.


On remarque que la propriété PO est un invariant affine en conséuqence de quoi,
toutes les restrictions de $f$ au hyperplan affine partagent la même 
non-linéaratité d'ordre $r$.






en utilisant  la valeur  $\rho(1,7)=56$ \cite{MYK,HOU},  et les valeurs obtenues récement $\rho(2,7)=40$ \cite{Wang19} 
et $\rho(3,7)=20$ \cite{GKLW}. L'objectif de cette note, est de monrer que les  majorations sont strictes. Pour cela,
nous devons établir que les fonctions de non linéarité quadratique optimale sont cubique ou quartique. Tout d'abord
en utilisant la classification de $\bst(3,4,7)$, nous trouvons 4 classes de fonctions de degré au plus 4
à distance 40 du Reed-Muller d'ordre 2 :

\begin{table}
\centering
\caption{\label{prop:4classes}
Les $4$ classes de fonctions de $\bst(3,4,7)$ de non linéarité quadratique 
 $40$ ; leurs degrés sont
$4,4,3,4$. Les distributions de $\Nl(f+q)$ sur les $2^{21}$ quadratiques homogènes $q$.
}
\begin{tabular}{lrrrr}
\toprule
classe & $40$ & $44$ & $48$ & $52$\\
\midrule
1 & 14\,336 & 290\,816 & 1\,598\,464 & 193\,536\\
2 & 14\,336 & 290\,816 & 1\,589\,248 & 202\,752\\
3 & 11\,264 & 307\,200 & 1\,768\,448 & 10\,240\\
4 & 11\,264 & 307\,200 & 1\,670\,144 & 108\,544\\
\bottomrule
\end{tabular}
\end{table}

\begin{result} Il  n'y a que 4 classes de fonctions de non linéarité quadratique 40,
ce sont les fonctions de degré 3 et 4 énumérées dans la table \ref{prop:4classes}. 
\end{result}
\begin{proof}
	La preuve numérique sera déroulée dans les sections qui suivent.
\end{proof}

\begin{proposition}
	Les fonctions hautement non quadratiques de 7 bits sont dans $\rmc(4,7)$ et
	nous avons $\rho(2,8) < 96$.
\end{proposition}
\begin{proof}
	Considérons $f$ une fonction de $\boole(8)$ de non linéarité 96. L'inégalité de Plotkin
	montre que $f = ( f_0 || f_1 )$ où $f_0$ et $f_1$ sont hautement non quadratique en 7
	variables. La distribution des non linéarité d'ordre 1 de ces fonctions montre sur le
	champs qu'une concaténation de deux fonctions hautement non quadratique est au plu égale
	à 94.
\end{proof}

\begin{proposition}
        Le rayon de recouvrement de \rmc(3,8) vérifie $\rho(3,8) < 60$.
\end{proposition}
\begin{proof}
        En effet, si $f$ une fonction de $\boole(8)$ de non linéarité  cubique 60 alors
	c'est la concaténation de deux fonctions de non linéarité cubique 20 dont les 
	meilleures 3-approximation sont de non linéarité quadratique 40. Il s'ensuit que
	$f\in\rmc(4,8)$ et on peut utiliser le résultat de \cite{GAO} qui affirme que
        le rayon de recouvrement relatif $\rho_4(3,8) =  56$  pour conclure.
\end{proof}


\section{Classification}
\label{sec:class}

Nous notons $\bst(s,t,m)$ l'espace des fonctions booléennes
de degré au plus $t$ et de valuation supérieure ou égale à $s$.
Deux fonctions Booléennes $f$ et $g$ de $\boole(m)$ sont dites $r$-équivalentes
quand il existe une transformation affine $\agl s \in\AGL(2,m)$ tel que $\deg( g+f\circ \agl s) \leq r$. 
Nous notons $\clst(s,t,m)$ un système de représentants sous 
cette action. Le nombre de classes est déterminé par la formule 
de Burnside  et $\clst(s, t, m)$ s'obtient par descente à partir de $\bst(s+1, t, m)$.  




Nous sommes partis des stabiliseurs des 179 classes  de $\bst(4,5,7)$ pour obtenir un système
de représentants de $\bst(3,5,7)$.   Le stabiliseur
de $f\in\bst(4,5,7)$ est le sous-groupe $S$ composés des  $\agl s$ tel que $\deg( f + f\circ \agl s ) \leq 3$.
Les orbites de $\bst(3,5,7)$ 3-equivalentes à $f$ ss'obtiennent par l'action 
$\bst(3,3,7)\ni c\mapsto  c\circ \agl s +( h + h\circ \agl s)$. Pour la réalisation, nous
avons distingué deux situations : suivant l'ordre de $S$.

La méthode de descente est décrite dans [VGPL],
une description des  $118\,140\,881\,980\approx 2^{36.78}$ 
classes est présentée dans la page du projet numérique \cite{PN}, 
sous la forme de 178 tableaux de $2^{29}$ bits, après  compression 
la donnée pèse moins d'un gigaoctet. 



\begin{result}
	Les quintiques de de $\boole(7)$ ne sont pas hautement quadratique.   
\end{result}
\begin{proof} On peut vérifier ce point en appliquant le décodage par symbole d'inforimation de Prange.
\end{proof}

\section{Structure des quintique}

Les mots de poids mininum du translaté de $f+\rmc(r,m)$ sont
les leaders associés à $f$. Convenons de dire que $f$ définit
un orphelin ou urcosets quand l'union des supports des leaders 
de $f$ est égal à $\fd(m)$. Un résultat fondamental de \cite{BRUALDI}
affirme que si $f$ définit un orphelin alors, pour tout $a\in\fd(m)$.
$\NL(r,f) = \NL(r,f+\delta_a)$.

Nous avons déterminé les classes dont la non-linéarité
quadratique est supérieure ou égale à 38 en combinant un 
filtre basé sur l'algorithme de
Prange avec une recherche explicite de leader.  Nous obtenons
51 fonctions : 1 cubique,  36 quartique et 14 quintique.  

L'ensemble des données de classification 
de $\bst(3,5,7)$ sont accessibles via la
page du project \cite{GL26}.

Il est bien connu que le rayon de recouvrement du Reed-Muller d'ordre
$m-2$ est égal à 2. Un résultat que nous pouvons préciser de la 
façon suivante :

\begin{lemma}
	Si $f$ est une fonction de degré $m-1$ alors il
	existe une fonction de degré exactement $m-2$ à
	distance 2 de $f$.  
\end{lemma}
\begin{proof}
Considérons $f$ de degré $m-1$, 
$$
f(x) \equiv \sum_{i} a_i \widehat{ x_i}  + \sum_{i<j} a_{ij} \widehat{x_{ij}}\mod {\rmc(m-3,m)} 
$$
Notons $a=(a_1,a_2,\ldots,a_m)$, pour tout $v\in\fd(m)$, 
la fonction $\delta_v( x ) + \delta_{a+v}( x ) + f( x )$ est nulle modulo $\rmc(m-1,m)$ 
et le coefficient  de degré $m-2$ est donné par :
$$
	 \sum_{i<j} a_i + a_j + a_i v_j + a_j v_i + a_i a_j  + a_{ij},
$$
Une fonction de degré 1 en $v$ qui n'est pas contante  
car au mois l'un des l'un des $a_i$ n'est pas nul.
\end{proof}

\begin{result}
Les 51 classes de non linéarité quadratique supérieure ou égale
à 38 fonctions définissent des orphelins du Reed-Muller d'ordre 2.
	Une modification ponctuelle de ces fonctions diminue 
	leur distance au Reed-Muller d'ordre 2.
\end{result}


\begin{proposition}
\label{prop:complet}
Si $f\in\boole(7)$ est de non linéarité quadratique $40$, 
alors $f$ est une fonction de degré au plus 4 équivalente 
à une des fonctions de la table \ref{TABLE}.
\end{proposition}
\begin{proof}
Nous savons que si $f$ n'est pas de degré $\leq 4$, elle ne peut pas être de degré 5, ni 7. 
Supposons que celle ci soit de degré 6.  Considérons une fonction $g=\delta_a+\delta_b+f$ 
de degré 5 à distance 2 de $f$ comme dans le Lemme \ref{MODIF}. La nonlinéarité de $g$ est
38, et elle définit un urcoset, il suit que $ \NL(2, g + \delta_a)  < 38$ 
et donc $\NL(2,f) < 40$, en contradiction avec notre hypothèse.
\end{proof}


\begin{coro}
	Le rayon de recouvrement de $\rmc(2, 8)$ vérifie $\rho(2,8) < 96$.
\end{coro}
\begin{proof}
\end{proof}


\begin{coro}
	Le rayon de recouvrement de $\rmc(3, 8)$ vérifie $\rho(3,8) < 60$.
\end{coro}
\begin{proof}
\end{proof}


% =====================================================================
\begin{thebibliography}{99}
\bibitem{Wang19} Q.~Wang, \emph{The covering radius of the Reed--Muller code $\rmc(2,7)$ is 40},
 Discrete Math. \textbf{342}(12), 111625, 2019.
\bibitem{GKLW} J.~Gao, H.~Kan, Y.~Li, Q.~Wang, \emph{The covering radius of the third-order
 Reed--Muller code $\rmc(3,7)$ is 20}, arXiv:2206.10881.
\bibitem{Gao26} J.~Gao, \emph{Numerical results and asymptotic lower bound on the covering
 radius of Reed--Muller codes $\rmc(2,11)$ and $\rmc(3,n)$}, IEICE Trans.
 Fundamentals \textbf{E109-A}(1), 35--46, 2026.
\bibitem{DMT22} R.~Dougherty, R.~D.~Mauldin, M.~Tiefenbruck, \emph{The covering radius of the
 Reed--Muller code $\rmc(m-4,m)$ in $\rmc(m-3,m)$}, IEEE Trans. Inform. Theory
 \textbf{68}(1), 560--571, 2022.
\bibitem{GL26} V.~Gillot, P.~Langevin, \emph{Classification of $\bst(3,5,7)$},\newline
\texttt{http://langevin.univ-tln.fr/project/B357/B357.html}
\bibitem{GL23} V.~Gillot, P.~Langevin, \emph{Covering radius of $\rmc(4,8)$},
 arXiv:2305.03493, 2023 ; projet agl8, \texttt{langevin.univ-tln.fr/project/agl8}.
\bibitem{LL05} P.~Langevin, G.~Leander, \emph{Classification of Boolean quartic forms in eight
 variables}, in \emph{Boolean Functions in Cryptology and Information Security}, NATO Sci.
 Peace Secur. Ser. D \textbf{18}, 139--147, 2005.
\bibitem{Hou96} X.-D.~Hou, \emph{$GL(m,2)$ acting on $R(r,m)/R(r-1,m)$}, Discrete Math.
 \textbf{149}, 99--122, 1996.
\bibitem{CL92} G.~Cohen, S.~Litsyn, \emph{On the covering radius of Reed--Muller codes},
 Discrete Math., 1992, doi:10.1016/0012-365X(92)90542-N.
\end{thebibliography}

\end{document}
